\documentclass[preview]{standalone}

\usepackage{amsmath}
\usepackage{amssymb}
\usepackage{stellar}
\usepackage{enumitem}
\usepackage{definitions}

\begin{document}

\id{linearalgebra-exercises-batch-14}
\genpage

\section{Exercises}

\begin{snippetexercise}{linear-algebra-batch-14-ex-1}{}
    Consider the endomorphism \(f \colon \realnumbers^3 \fromto \realnumbers^3\) represented,
    with respect to the standard basis, by the matrix
    \[A = \begin{pmatrix} 14 & -13 & 8 \\ -13 & 14 & 8 \\ 8 & 8 & -7 \end{pmatrix}.\]
    Find an orthonormal basis of \(\realnumbers^3\) consisting of eigenvectors of \(f\).
\end{snippetexercise}

\begin{snippetsolution}{linear-algebra-batch-14-ex-1-sol}{}
    \todo
\end{snippetsolution}

\begin{snippetexercise}{linear-algebra-batch-14-ex-2}{}
    Consider \(\mathcal{P}_2(\realnumbers)\) with the inner product \[\langle p, q \rangle \triangleq \int_0^1 pq\] and the operator \(f \in \operatorname{End}(\mathcal{P}_2(\realnumbers))\), \(f(a_0 + a_1x + a_2x^2) = a_1x\).
    \begin{enumerate}[label=\Alph*]
        \item Prove that \(f\) is not self-adjoint.
        \item Prove that the matrix associated with \(f\) in the basis \(\{1, x, x^2\}\) is \(D = \operatorname{diag}(0, 1, 0)\).
        \item \(D\) is symmetric even though \(f\) is not self-adjoint. Is there a contradiction?
    \end{enumerate}
\end{snippetexercise}

\begin{snippetsolution}{linear-algebra-batch-14-ex-2-sol}{}
    \todo
\end{snippetsolution}

\begin{snippetexercise}{linear-algebra-batch-14-ex-3}{}
    The \emph{orthogonal group} and the
    \emph{unitary group} are respectively
    \[\operatorname{O}(n) = \{A \in \operatorname{GL}(n; \realnumbers) \suchthat AA^t = I_n\} \quad \text{and} \quad \operatorname{U}(n) = \{A \in \operatorname{GL}(n; \complexnumbers) \suchthat A\bar{A}^t = I_n\}.\]
    \begin{enumerate}[label=\Alph*]
        \item Prove that \(\operatorname{O}(n)\) and \(\operatorname{U}(n)\) are groups.
        \item Prove that \(A \in \operatorname{O}(n)\) (resp. \(A \in \operatorname{U}(n)\)) if
        and only if the columns of \(A\) form an orthonormal basis of
        \(\realnumbers^n\) (resp. \(\complexnumbers^n\)) with respect to the standard inner product.
        \item Prove that if \(A \in \operatorname{O}(n)\), then \(\det(A) \in \{1, -1\}\); is the converse true?
        \item Prove that if \(A \in \operatorname{U}(n)\), then
        \(\det(A) \in S^1 = \{z \in \complexnumbers \suchthat |z| = 1\}\);
        is the converse true?
    \end{enumerate}
\end{snippetexercise}

\begin{snippetsolution}{linear-algebra-batch-14-ex-3-sol}{}
    \todo
\end{snippetsolution}

\begin{snippetexercise}{linear-algebra-batch-14-ex-4}{}
    In this exercise we discuss some isometries of \((\realnumbers^n, \cdot)\).
    \begin{enumerate}[label=\Alph*]
        \item Prove that if \(\lambda \in \realnumbers\) is an eigenvalue of
        \(A \in \operatorname{O}(n)\), then \(\lambda = \pm 1\).
        \item Prove that if \(n\) is odd, then
        \(A \in \operatorname{O}(n)\) has at least one real eigenvalue.
        \item Let \(f \colon \realnumbers^3 \fromto \realnumbers^3\) be an isometry with
        \(\det f = 1\) and let
        \(A = \mathcal{M}(f, \mathcal{C}an) \in \operatorname{O}(3)\).
        Prove that there exists an orthonormal basis \(\mathcal{B}\) of \(\realnumbers^3\) such
        that \(A = P^{-1}BP\), with
        \(P = \mathcal{M}(\operatorname{Id}_{\realnumbers^3}, \mathcal{C}an, \mathcal{B})\) and
        \[B = \begin{pmatrix} 1 & 0 & 0 \\ 0 & \cos \theta & -\sin \theta \\ 0 & \sin \theta & \cos \theta \end{pmatrix}\] for some
        \(\theta \in \realnumbers\).
        In other words, \(f\) is a rotation by an angle \(\theta\) in the plane orthogonal
        to the first vector of \(\mathcal{B}\), which is the axis of rotation.
        \item Calculate the axis and angle of the rotation represented, with respect to the
        standard basis, by the matrix
        \[A = \begin{pmatrix} 0 & 0 & 1 \\ 1 & 0 & 0 \\ 0 & 1 & 0 \end{pmatrix}.\]
    \end{enumerate}
\end{snippetexercise}

\begin{snippetsolution}{linear-algebra-batch-14-ex-4-sol}{}
    \todo
\end{snippetsolution}

\begin{snippetexercise}{linear-algebra-batch-14-ex-5}{}
    Determine the sign of the following quadratic forms:
    \begin{enumerate}
        \item \(q \colon \realnumbers^3 \fromto \realnumbers^3\) given by \(q(x,y,z) \triangleq -2x^2 + 2xy + 2xz - 2y^2 + 2yz - 2z^2\);
        \item \(f \colon \mathcal P(\realnumbers) \fromto \realnumbers\) given by \(f(p) \triangleq p(0)^2 - p(1)^2 + p(-1)^2\);
        \item \(g \colon \matrices_{2\times 2}(\realnumbers) \fromto \realnumbers\) given by \(g(A) \triangleq \trace(A^2)\).
    \end{enumerate}
\end{snippetexercise}

\begin{snippetsolution}{linear-algebra-batch-14-ex-5-sol}{}
    \todo
\end{snippetsolution}

\begin{snippetexercise}{linear-algebra-batch-14-ex-6}{}
    Show that the \matrix \(A = \begin{pmatrix}
        a & b \\ c & d
    \end{pmatrix} \in \matrices_{2\times 2}(\realnumbers)\) is:
    \begin{itemize}
        \item positive definite (respectively negative) if \(a>0\) (respectively \(a<0\)) and \(ac-b^2 >0\);
        \item positive semidefinite (respectively negative) if \(a\geq 0\) (respectively \(a \leq 0\)) and \(ac-b^2 = 0\);
        \item indefinite if \(ac-b^2 < 0\).
    \end{itemize}
\end{snippetexercise}

\begin{snippetsolution}{linear-algebra-batch-14-ex-6-sol}{}
    \todo
\end{snippetsolution}

\begin{snippetexercise}{linear-algebra-batch-14-ex-7}{}
    Given the symmetrical \matrix
    \[
        B = \begin{pmatrix}
            2 & -1 & 0 \\ -1 & 2 & -1 \\ 0 & -1 & 2
        \end{pmatrix} \in \matrices_{3 \times 3}(\realnumbers)
    \]
    Show that, for column vectors
    \(x,y \in \realnumbers^3\), \(\langle x,y\rangle \triangleq x^t B y\) defines
    an inner product on \(\realnumbers^3\).
\end{snippetexercise}

\begin{snippetsolution}{linear-algebra-batch-14-ex-7-sol}{}
    \todo
\end{snippetsolution}

\end{document}